\uuid{Q6ys}
\titre{Matrice d'une application linéaire/Polynômes}
\auteur{hanani ; mzari}
\organisation{exo7}
\licence{CC BY-NC-SA 4.0}
\datecreate{2020-01-20}
\niveau{L1}
\module{Algèbre}
\chapitre{Matrice}
\sousChapitre{Matrice et application linéaire}
\difficulte{3}
\theme{Matrice d'une application linéaire, Polynômes}
\exo7qcmid{lille-2.matrice-application.016}

\contenu{
\begin{question}
\qclassification{108.03}
\qtags{motcle=Matrice d'une application linéaire/Polynômes}
On considère $\Rr_2[X]$ l'ensemble des polynômes à coefficients réels de degré $\le 2$ et  $\Rr^3$ munis de leurs bases canoniques notées respectivement ${\cal {B}}_1$ et ${\cal {B}}$. Soit $f$ l'application linéaire définie par  :
$$\begin{array}{rccc}f:&\Rr_2[X] &\to& \Rr^3\\
& P&\to &\big(P(0),P(1),P(-1)\big).  \end{array}$$
On considère la base ${\cal {B}}_2 = \{P_1,P_2,P_3\}$, où $ P_1=1, P_2=1+X, P_3=1+X^2$. Quelles sont les assertions vraies ?
\begin{answers}
\bad{La matrice de $f$ de la base $ {\cal {B}}_1$ à la base  ${\cal {B}}$ est : $
\left(\begin{array}{rcc}
1&1&1\\
0&1&-1\\
0&1&1\\
\end{array}\right).$}
\good{La matrice de $f$ de la base $ {\cal {B}}_1$ à la base  ${\cal {B}}$  est : $
\left(\begin{array}{rcc}
1&0&0\\
1&1&1\\
1&-1&1\\
\end{array}\right).$}
\good{La matrice de $f$ de la base $ {\cal {B}}_2$ à la base  ${\cal {B}}$  est : $
\left(\begin{array}{rcc}
1&1&1\\
1&2&2\\
1&0&2\\
\end{array}\right).$}
\bad{La matrice de $f$ de la base $ {\cal {B}}_2$ à la base  ${\cal {B}}$  est : $
\left(\begin{array}{rcc}
1&1&1\\
1&2&0\\
1&3&1\\
\end{array}\right).$}
\end{answers}
\vskip2mm
\begin{explanations} La matrice de $f$  d'une base $ {\cal {B}}=(u_j)$ dans une  base $ {\cal {B}}'=(v_i)$ est la matrice  dont la $j$ième colonne est constituée des coordonnées de  $f(u_j)$ dans la base $ {\cal {B}}'$. La matrice de $f$ de la base $ {\cal {B}}_1$ à la base  ${\cal {B}}$  est : $
\left(\begin{array}{rcc}
1&0&0\\
1&1&1\\
1&-1&1\\
\end{array}\right).$
\vskip0mm
La matrice de $f$ de la base $ {\cal {B}}_2$ à la base  ${\cal {B}}$  est : $
 \left(\begin{array}{rcc}
1&1&1\\
1&2&2\\
1&0&2\\
\end{array}\right).$
\end{explanations}
\end{question}
}
